\documentclass{article}
\usepackage{geometry}
\usepackage{amsmath}
\usepackage{listings}
\usepackage{xcolor}
\usepackage{hyperref}

\geometry{a4paper, margin=1in}

% --- lstlisting Configuration ---
\definecolor{codegreen}{rgb}{0,0.6,0}
\definecolor{codegray}{rgb}{0.5,0.5,0.5}
\definecolor{codepurple}{rgb}{0.58,0,0.82}
\definecolor{backcolour}{rgb}{0.95,0.95,0.92}

\lstdefinestyle{mystyle}{
    backgroundcolor=\color{backcolour},
    commentstyle=\color{codegreen},
    keywordstyle=\color{magenta},
    numberstyle=\tiny\color{codegray},
    stringstyle=\color{codepurple},
    basicstyle=\ttfamily\footnotesize,
    breakatwhitespace=false,
    breaklines=true,
    captionpos=b,
    keepspaces=true,
    numbers=left,
    numbersep=5pt,
    showspaces=false,
    showstringspaces=false,
    showtabs=false,
    tabsize=2
}
\lstset{style=mystyle}
% --- End lstlisting Configuration ---

\title{Pyomo Implementation Pipeline for Huawei TechArena 2025}
\author{Gen's BESS Optimization Team}
\date{\today}

\begin{document}
\maketitle

This document outlines the step-by-step process for implementing the BESS optimization model using Python and Pyomo.

\section{Project Setup}

\subsection{Directory Structure}
A logical project structure is crucial for maintainability.
\begin{lstlisting}[language=bash, caption={Recommended Directory Structure}]
.
|-- input/
|   `-- TechArena2025_data.jsonl
|-- output/
|-- main.py
`-- requirements.txt
\end{lstlisting}

\subsection{Environment and Dependencies}
A dedicated virtual environment prevents dependency conflicts. The \texttt{requirements.txt} file should specify all necessary packages.
\begin{lstlisting}[caption={\texttt{requirements.txt}}]
pandas
numpy
pyomo
openpyxl
# A solver like CBC or GLPK must be installed and accessible in the system PATH.
\end{lstlisting}

\section{Data Loading and Preprocessing (\texttt{main.py})}
This phase transforms raw data into a model-ready format.

\subsection{Load Raw Data}
\begin{itemize}
    \item Read the \texttt{TechArena2025\_data\_tidy.jsonl} file line by line into a pandas DataFrame.
\end{itemize}

\subsection{Clean and Structure Data}
\begin{itemize}
    \item Convert the \texttt{timestamp} column to datetime objects and set it as the DataFrame index.
    \item Pivot the DataFrame to create a multi-level column structure: \texttt{(country, source, direction)}.
\end{itemize}

\subsection{Resample and Align Time Series}
\begin{itemize}
    \item Resample all time series to a consistent 15-minute frequency (\texttt{15T}).
    \item Use forward-fill (\texttt{ffill}) to propagate prices within their validity periods (e.g., a 4-hour ancillary price applies to all 15-min intervals within that block).
    \item Create helper columns for 'day' and '4-hour block' identifiers to facilitate constraint mapping.
\end{itemize}

\section{Main Loop Structure (\texttt{main.py})}
The core logic iterates through every required scenario (country $\times$ C-rate $\times$ cycle limit).
\begin{lstlisting}[language=Python, caption={Pseudocode for the Main Scenario Loop}]
countries = ['DE', 'AT', 'CH', 'HU', 'CZ']
c_rates = [0.25, 0.33, 0.50]
daily_cycles = [1.0, 1.5, 2.0]

all_scenario_results = {}

for country in countries:
    for c_rate in c_rates:
        for n_cycles in daily_cycles:
            # Get the preprocessed data for the current country
            country_data =...

            # Define BESS parameters for this specific scenario
            bess_params = {
                'E_nom': 4472.0,  # kWh [1]
                'P_max_config': 4472.0 * c_rate,  # kW [1]
                'N_cycles': n_cycles,
                'eta_ch': 0.90**0.5,  # Round-trip efficiency is 90%
                'eta_dis': 0.90**0.5,
                'SOC_min': 0.0,
                'SOC_max': 1.0,
                'delta_t': 0.25,  # hours
                'delta_b': 4.0, # hours
                'min_bid_fcr': 1.0, # MW
                'min_bid_afrr': 1.0 # MW
            }
            
            # Apply country-specific rules (CRITICAL)
            if country == 'CH':
                # Swiss aFRR market has a 5 MW minimum bid, BESS max power is ~2.2 MW
                # Therefore, the BESS cannot participate in the Swiss aFRR market.
                bess_params['can_bid_afrr'] = False
            else:
                bess_params['can_bid_afrr'] = True

            # Call the optimization function
            # This should implement the rolling horizon strategy (see Section 6)
            solution = solve_bess_optimization(country_data, bess_params)

            # Store the results
            all_scenario_results[(country, c_rate, n_cycles)] = solution

# 5. Post-process all_results and generate output files
\end{lstlisting}

\section{Pyomo Model Implementation (\texttt{run\_optimization} function)}
This function encapsulates the mathematical model definition and solver execution. Follow the official Pyomo documentation for detailed syntax and capabilities: \url{https://pyomo.readthedocs.io/en/stable/}.

\subsection{Initialize Model and Define Sets}
\begin{lstlisting}[language=Python]
import pyomo.environ as pyo

def run_optimization(data, params):
    model = pyo.ConcreteModel(name="BESS_Optimization")
    
    # Define Sets from data indices
    model.T = pyo.Set(initialize=range(len(data.index)))
    model.B = pyo.Set(initialize=data['block_id'].unique())
    model.D = pyo.Set(initialize=data['day_id'].unique())
\end{lstlisting}

\subsection{Define Parameters and Variables}
\begin{itemize}
    \item Create \texttt{pyo.Param} objects for all input data (prices, efficiencies, etc.).
    \item Create \texttt{pyo.Var} objects for all decision variables, specifying domains (e.g., \texttt{pyo.NonNegativeReals}, \texttt{pyo.Binary}).
\end{itemize}

\subsection{Define Objective Function and Constraints}
Each mathematical equation is translated into a Python function.
\begin{lstlisting}[language=Python, caption={Example Constraint Implementation}]
# --- Objective Function ---
def objective_rule(model):
    # Return objective expression from the refined mathematical model
    return sum(...)
model.objective = pyo.Objective(rule=objective_rule, sense=pyo.maximize)

# --- SOC Dynamics Constraint ---
def soc_dynamics_rule(model, t):
    if t == model.T.first():
        return model.e_soc[t] == params['E_soc_initial'] + ...
    return model.e_soc[t] == model.e_soc[t-1] + ...
model.soc_dynamics = pyo.Constraint(model.T, rule=soc_dynamics_rule)

# --- Daily Cycle Constraint ---
def daily_cycle_rule(model, d):
    # Get all time steps 't' for day 'd'
    time_steps_in_day = [t for t in model.T if data.loc[t, 'day_id'] == d]
    return sum(model.p_ch[t] * params['dt'] for t in time_steps_in_day) <= \
           model.N_cycles * model.E_nom
model.daily_cycle_limit = pyo.Constraint(model.D, rule=daily_cycle_rule)
\end{lstlisting}

\subsection{Solve and Extract Results}
\begin{itemize}
    \item Instantiate a solver (e.g., CBC, GLPK) and call the \texttt{solve()} method.
    \item Extract the optimal values from the model's variables.
    \item Return a structured object (e.g., a dictionary) containing the key results for the scenario.
\end{itemize}

\section{Post-Processing and Output Generation (\texttt{main.py})}
The final step aggregates results and formats them according to submission guidelines.

\subsection{Aggregate Results and Calculate Metrics}
\begin{itemize}
    \item Convert the list of result dictionaries into pandas DataFrames.
    \item Calculate final metrics like Levelized Return on Investment (ROI) over a 10-year horizon.
\end{itemize}

\subsection{Write to Excel}
As per the FaQ, create three Excel files, each containing separate sheets for each country.
\begin{itemize}
    \item \texttt{TechArena\_Phase1\_Configuration.xlsx}
    \item \texttt{TechArena\_Phase1\_Investment.xlsx}
    \item \texttt{TechArena\_Phase1\_Operation.xlsx}
\end{itemize}
\begin{lstlisting}[language=Python, caption={Pseudocode for Writing Excel Output}]
# Pseudocode for writing one of the output files
with pd.ExcelWriter('output/TechArena_Phase1_Configuration.xlsx') as writer:
    for country in countries:
        df_config_country = df_all_configs[df_all_configs['Country'] == country]
        df_config_country.to_excel(writer, sheet_name=country, index=False)
# Repeat for Investment and Operation files
\end{lstlisting}

\end{document}
